% chapters/chapter3/3_1_intro.tex
% !TeX root=../../main.tex

\section{چارچوب ارزیابی و سناریوهای شبیه‌سازی}

در این فصل، معماری‌های کنترلی توسعه‌یافته بر روی فیزیک غیرخطی سیستم در بستر نرم‌افزار \lr{Simscape Multibody} استقرار یافته و مورد ارزیابی قرار می‌گیرند. هدف از این شبیه‌سازی‌ها، بررسی گام‌به‌گام کارایی، محدودیت‌ها و مصالحه‌های مهندسی (\lr{Trade-offs}) در هدایت کالسکه، مدیریت دینامیک زیرفعال آونگ و نحوه برخورد با محدودیت‌های فیزیکی عملگرهاست.

برای تحلیل دقیق و مقایسه‌ای عملکرد، چهار ساختار کنترلی متفاوت ارزیابی می‌شوند:
\begin{enumerate}
    \item \textbf{\lr{LQR-Step}:} کنترل‌کننده خطی در مواجهه مستقیم با خطای موقعیت ناگهانی (ورودی پله).
    \item \textbf{\lr{LQR-Trajectory}:} کنترل‌کننده خطی ادغام‌شده با فیدفوروارد تحلیلی و برنامه‌ریز مسیر سینماتیکی.
    \item \textbf{\lr{NMPC-Step}:} کنترل‌کننده پیش‌بین غیرخطی در مواجهه مستقیم با خطای موقعیت.
    \item \textbf{\lr{NMPC-Trajectory}:} معماری پیش‌بین غیرخطی تغذیه‌شده توسط افق پیش‌بینی مسیر مرجع.
\end{enumerate}

به منظور به چالش کشیدن پایداری دینامیکی و نمایان ساختن تفاوت رفتار این کنترل‌کننده‌ها در شرایط عملیاتی مختلف، تمامی ارزیابی‌ها تحت دو سناریوی جابه‌جایی استاندارد انجام می‌پذیرد:
\begin{itemize}
    \item \textbf{مانور کوتاه‌برد ($0.5$ متر):} طراحی‌شده برای سنجش رفتار سیستم در مجاورت نقطه تعادل و بررسی اعتبار تقریب‌های خطی و ردیابی گذرا.
    \item \textbf{مانور بلندبرد ($5$ متر):} طراحی‌شده برای ارزیابی تاب‌آوری سیستم در برابر خطاهای بزرگ، مدیریت قیود فیزیکی عملگرها و قابلیت میراسازی نوسانات پس از شتاب‌گیری‌های طولانی.
\end{itemize}

در تحلیل‌های آتی، عملکرد هر ساختار بر مبنای معیارهای مشخصی نظیر دقت ردیابی موقعیت، توانایی محدودسازی زاویه نوسان در کران مجاز ($\pm15^\circ$)، مهار نیروی درخواستی درون محدوده ظرفیت فیزیکی سیستم ($50\,\mathrm{N}$)، سرکوب کامل نوسان پسماند (\lr{Residual Sway}) و کارآمدی استفاده از درجه آزادی افزوده طول کابل به عنوان یک مکانیزم میرایی فعال ارزیابی خواهد شد.

روند استدلال در این فصل به صورت علی تدوین شده است؛ بدین معنا که ناتوانی‌های بنیادی رؤیت‌شده در پاسخ هر کنترل‌کننده، ضرورت گذار و طراحی معماری پیچیده‌تر بعدی را توجیه و اثبات می‌نماید. نهایتاً در بخش‌های پایانی، با تکیه بر تحلیل‌های کمی، تفسیر فیزیکی و مقایسه جامع دستاوردها و هزینه‌های محاسباتی، جمع‌بندی واقع‌بینانه‌ای برای پیاده‌سازی این الگوریتم‌ها در کاربردهای صنعتی ارائه می‌گردد.